\section{The zeta function of an abelian variety}
\begin{definition}[Frobenius and zeta function]\label{def:milne-frobenius-zeta}\dcref{MI:ch2:1}
Write $\overline{\F_q}$ for an algebraic closure of $\F_q$ and
$\F_{q^m}$ for its unique subfield with $q^m$ elements. Let $V/\F_q$ be a
variety. The Frobenius map $\pi_V:V\to V$ is the identity on the underlying
topological space and sends a local function $f$ to $f^q$. For
$V=\PP^n$, it is induced by $X_i\mapsto X_i^q$ and sends
$(x_0:\cdots:x_n)$ to $(x_0^q:\cdots:x_n^q)$. Every morphism over $\F_q$
commutes with Frobenius, and $V(\F_{q^m})$ is the fixed-point set of
$\pi_V^m$ on $V(\overline{\F_q})$. For $N_m=|V(\F_{q^m})|$, its zeta function is
\[
 Z(V,t)=\exp\left(\sum_{m\geq 1}N_m\frac{t^m}{m}\right).
\]
\end{definition}

\begin{theorem}[Weil bounds]\label{thm:milne-zeta-weil-bounds}\dcref{MI:ch2:1.1}
Let $A/\F_q$ be an abelian variety of dimension $g$, let $\pi$ be its
Frobenius, and write its characteristic polynomial as
$P_\pi(X)=\prod_{i=1}^{2g}(X-a_i)$. If $N_m=|A(\F_{q^m})|$, then
$N_m=\prod_i(1-a_i^m)$ and $|a_i|=q^{1/2}$. Moreover,
\[
 |N_m-q^{mg}|\leq 2gq^{m(g-1/2)}+(2^{2g}-2g-1)q^{m(g-1)}.
\]
\end{theorem}

\begin{proof}[Source proof]

We first deduce the inequality from the preceding statements.  Take $m=1$
in (a) and expand the product.  The term containing all $2g$ eigenvalues is
an integer, in fact the positive integer $P_\pi(0)=\deg(\pi)$; by (b) its
absolute value is $q^g$, and hence it is $q^g$.  Each of the $2g$ terms in
which one eigenvalue is omitted has absolute value $q^{g-1/2}$, while the
remaining $2^{2g}-2g-1$ terms have absolute value at most $q^{g-1}$.
Replacing $\pi$ by $\pi^m$ gives the displayed inequality.

We now prove (a).  First take $m=1$.  The kernel of
$\pi-\operatorname{id}:A(\F)\to A(\F)$ is $A(\F_q)$.  The differential of Frobenius at
the origin is zero.  Indeed, after replacing an affine neighbourhood of the
origin by an affine space with the origin at $0$, Frobenius is given by
$Y_i=X_i^q$, whose differential is zero in characteristic $p$.  Therefore
$d(\pi-\operatorname{id})_0=-1$, so $\pi-\operatorname{id}$ is etale at the origin and, by homogeneity,
at every point.  Every point of its kernel consequently occurs with
multiplicity one, and
\[
 |A(\F_q)|=\deg(\pi-\operatorname{id})=P_\pi(1).
\]
The same argument applied to $\pi^m$ gives
$|A(\F_{q^m})|=P_{\pi^m}(1)$.  The eigenvalues of $\pi^m$ on $V_\ell A$
are $a_1^m,\ldots,a_{2g}^m$, so
\[
 P_{\pi^m}(X)=\prod_i(X-a_i^m),\qquad
 P_{\pi^m}(1)=\prod_i(1-a_i^m),
\]
as required.  Part (b) follows from the next two lemmas.

\end{proof}

\begin{lemma}[Rosati--Frobenius identity]\label{lem:milne-rosati-frobenius}\dcref{MI:ch2:1.2}
Let $\dagger$ be the Rosati involution on $\End(A)\otimes\Q$ defined by a
polarization of $A/\F_q$. Then
$\pi_A^\dagger\circ\pi_A=q\,\operatorname{id}_A$.
\end{lemma}

\begin{proof}[Source proof]
Let $D$ be the ample divisor defining the polarization, so that
$\lambda(a)=[D_a-D]$. It is enough to prove
$\pi^\vee\circ\lambda\circ\pi=q\lambda$. If $D'$ is a divisor and
$D'=\operatorname{div}(f)$ near $\pi(P)$, then
$\pi^*D'=\operatorname{div}(f\circ\pi)=qD'$ near $P$; hence
$\pi^*D=qD$. Also, for every homomorphism $\alpha:A\to A$ and point $a$,
$\alpha\circ t_a=t_{\alpha(a)}\circ\alpha$.

For $a\in A(\F)$, we have
\[
\begin{aligned}
 (\pi^\vee\circ\lambda\circ\pi)(a)
 &=\pi^*[D_{\pi(a)}-D] \\
 &=[(t_{\pi(a)}\circ\pi)^*D-\pi^*D] \\
 &=[(\pi\circ t_a)^*D-\pi^*D] \\
 &=[t_a^*(qD)-qD]=q\lambda(a).
\end{aligned}
\]
Thus $\pi^\vee\circ\lambda\circ\pi=q\lambda$, as required.
\end{proof}

\begin{lemma}[Rosati eigenvalue lemma]\label{lem:milne-rosati-eigenvalues}\dcref{MI:ch2:1.3}
Let $A$ be an abelian variety over a field $k$, and let $\alpha$ be an
element of $\End(A)\otimes\Q$ such that $\alpha^\dagger\alpha$ is an integer
$r$, where $\dagger$ is the Rosati involution defined by a polarization.  For
any root $a$ of $P_\alpha$ in $\C$, one has $|a|^2=r$.
\end{lemma}

\begin{proof}[Source proof]

The ring $\Q[\alpha]$ is a commutative finite-dimensional $\Q$-algebra and
therefore an Artin ring.  Let $\mathfrak m_1,\ldots,\mathfrak m_n$ be its
maximal ideals.  The quotient by their intersection is a product of fields,
\[
 \Q[\alpha]/\bigcap_i\mathfrak m_i=K_1\times\cdots\times K_n,
 \qquad K_i=\Q[\alpha]/\mathfrak m_i.
\]
We first show that $\bigcap_i\mathfrak m_i=0$.  The ring is stable under
the Rosati involution.  If $a\ne0$ in $\Q[\alpha]$, then
$b=a^\dagger a\ne0$, because $\operatorname{Tr}(a^\dagger a)>0$.  Since
$b^\dagger=b$, we have
$\operatorname{Tr}(b^2)=\operatorname{Tr}(b^\dagger b)>0$, so $b^2\ne0$;
iterating gives $b^{2^j}\ne0$ for every $j$.  Thus $b$, and hence $a$, is
not nilpotent.

Every automorphism $\tau$ of $\Q[\alpha]$ permutes the maximal ideals and
therefore the factors: there is a permutation $\sigma$ and isomorphisms
$\tau_i:K_i\to K_{\sigma(i)}$ such that
$\tau(a_1,\ldots,a_n)=(b_1,\ldots,b_n)$ with
$b_{\sigma(i)}=\tau_i(a_i)$.  For $\tau=\dagger$, the permutation must be
trivial, since otherwise positivity would fail on an element such as
$(a_1,0,\ldots,0)$.  Hence $\dagger$ preserves every factor and is positive
definite on each.

The involution extends by linearity (equivalently, by continuity) to a
positive-definite involution on $\Q[\alpha]\otimes_\Q\R$.  The same argument
shows that this is a product of fields, with each factor isomorphic to $\R$
or $\C$, and that the involution preserves every factor.  It is the identity
on a real factor.  On a complex factor it cannot be the identity, since that
is not positive definite, so it is complex conjugation.  Consequently, for
every homomorphism $\rho:\Q[\alpha]\to\C$,
\[
 \rho(\alpha^\dagger)=\overline{\rho(\alpha)}.
\]
Applying $\rho$ to $\alpha^\dagger\alpha=r$ gives
$|\rho(\alpha)|^2=r$.  The values $\rho(\alpha)$ are the roots of the
minimum polynomial of $\alpha$, and the roots of $P_\alpha$ are among these
values by the characteristic-polynomial comparison recalled in Chapter I.
This proves the lemma.

\end{proof}

\begin{corollary}[Zeta function]\label{cor:milne-zeta-function}\dcref{MI:ch2:1.5}
For $0\leq r\leq 2g$, let $P_r(t)=\prod_{i_1<\cdots<i_r}(1-a_{i_1}\cdots a_{i_r}t)$, with $P_0(t)=1$. Then
\[
 Z(A,t)=\frac{P_1(t)P_3(t)\cdots P_{2g-1}(t)}
 {P_0(t)P_2(t)\cdots P_{2g}(t)}.
\]
\end{corollary}

\begin{proof}
Take logarithms and use $\log(1-t)=-\sum_{m\ge1}t^m/m$ together with
Theorem~\ref{thm:milne-zeta-weil-bounds}; the coefficient of $t^m/m$ is
the alternating sum of the power sums of the $a_i^m$, which is $N_m$.
\end{proof}


\section{Abelian varieties over finite fields}
\begin{definition}[Category up to isogeny]\label{def:milne-isogeny-category}\dcref{MI:ch2:2}
For a field $k$, let $\operatorname{Isab}(k)$ be the $\Q$-linear category whose objects are
abelian varieties over $k$ and whose morphisms are
$\Hom_{\operatorname{Isab}(k)}(A,B)=\Hom(A,B)\otimes\Q$. Two objects are
isomorphic in $\operatorname{Isab}(k)$ exactly when they are isogenous. It is
additive, and every object is a direct sum of finitely many simple objects.
\end{definition}

\begin{lemma}[Frobenius and isogeny classes]\label{lem:milne-frobenius-isogeny}\dcref{MI:ch2:2.0}
Let $A$ be an abelian variety over $\F_q$ and let $\pi_A$ be its Frobenius.
Then $\pi_A$ commutes with every endomorphism of $A$. If $A$ is simple,
$\End^0(A)$ is a division algebra and $\Q[\pi_A]$ is a field. An isogeny
$A\to B$ between simple abelian varieties induces an isomorphism
$\End^0(A)\to\End^0(B)$ carrying $\pi_A$ to $\pi_B$.
\end{lemma}

\begin{proof}[Source proof]
Every morphism defined over $\F_q$ commutes with Frobenius. For simple $A$,
every nonzero endomorphism is an isogeny, so it is invertible in
$\End^0(A)$. If $A\to B$ is an isogeny, a quasi-inverse conjugates the two
endomorphism algebras and carries Frobenius to Frobenius.
\end{proof}

\begin{definition}[Weil integers and conjugacy]\label{def:milne-weil-integer}\dcref{MI:ch2:2}
A Weil $q$-integer is an algebraic integer $\pi$ such that, for every
embedding $\sigma:\Q[\pi]\hookrightarrow\C$,
$|\sigma(\pi)|=q^{1/2}$. Let $W(q)$ be the set of Weil $q$-integers in
$\overline{\Q}$. Two elements $\pi$ and $\pi'$ of $W(q)$ are conjugate if any
of the following equivalent conditions holds: they have the same minimum
polynomial over $\Q$; there is an isomorphism $\Q[\pi]\to\Q[\pi']$ carrying
$\pi$ to $\pi'$; or they lie in the same orbit under
$\operatorname{Gal}(\overline{\Q}/\Q)$.
\end{definition}

\begin{definition}[Central simple algebras and the Brauer group]\label{def:milne-central-simple}\dcref{MI:ch2:2}
A central simple $k$-algebra is a finite-dimensional simple $k$-algebra with
centre $k$; if it is a division algebra it is central division. The Brauer
group $\Br(k)$ consists of their similarity classes, with multiplication
induced by tensor product and inverse induced by the opposite algebra.
\end{definition}

\begin{theorem}[Honda--Tate classification]\label{thm:milne-honda-tate}\dcref{MI:ch2:2.1}
The map $A\mapsto\pi_A$ gives a bijection between isogeny classes of simple abelian varieties over $\F_q$ and conjugacy classes of Weil $q$-integers.
\end{theorem}

\begin{proof}[Source proof]

For injectivity, Tate's isogeny theorem gives
\[
 \Hom(A,B)\otimes\Q_\ell\simeq
 \Hom_{\operatorname{Gal}(\overline{\F}_q/\F_q)}(V_\ell A,V_\ell B).
\]
Frobenius acts semisimply, and the right side is nonzero when the Frobenius
polynomials have a common root. For simple $A$ and $B$, a nonzero homomorphism
is an isogeny, so equal conjugacy classes give equal isogeny classes.
More explicitly, if
$P_{\pi_A}(X)=\prod_i(X-a_i)$ and
$P_{\pi_B}(X)=\prod_j(X-b_j)$, then semisimplicity gives
\[
 \dim_{\Q_\ell}\Hom_{\operatorname{Gal}}(V_\ell A,V_\ell B)
 =\#\{(i,j):a_i=b_j\}.
\]
For surjectivity, Honda constructs the required varieties by starting with
complex-multiplication varieties over $\C$, choosing models over number fields,
and reducing at a suitable prime. Proposition~\ref{prop:milne-cm-reduction}
and Theorem~\ref{thm:milne-shimura-taniyama} identify the resulting Weil
integer; the remaining number-theoretic argument is Honda's theorem.

\end{proof}

\begin{lemma}[Tensor products of central simple algebras]\label{lem:milne-central-simple-tensor}\dcref{MI:ch2:2.2}
The tensor product over $k$ of two central simple $k$-algebras is central simple.
\end{lemma}

\begin{proof}[Source proof]

This is the standard central-simple tensor theorem (Milne cites CFT, IV, 2.8).

\end{proof}

\begin{proposition}[Wedderburn]\label{prop:milne-wedderburn}\dcref{MI:ch2:2.3}
Every central simple algebra is $M_r(D)$ for a central division algebra $D$, with $r$ uniquely determined and $D$ unique up to isomorphism.
\end{proposition}

\begin{proof}[Source proof]

This is Wedderburn's structure theorem (Milne cites CFT, IV, 1.15 and 1.20).

\end{proof}

\begin{theorem}[Local Brauer invariants]\label{thm:milne-local-brauer}\dcref{MI:ch2:2.4}
For every local field $k$, $\operatorname{inv}:\Br(k)\hookrightarrow\Q/\Z$ is canonical and injective; it is an isomorphism for nonarchimedean $k$, has image $\frac12\Z/\Z$ for $\R$, and $\Br(\C)=0$.
\end{theorem}

\begin{proof}[Source proof]

This is the local invariant theorem for Brauer groups (Milne cites CFT, III,
2.1 and IV, 4.3).

\end{proof}

\begin{theorem}[Global Brauer sequence]\label{thm:milne-global-brauer}\dcref{MI:ch2:2.6}
For a number field $k$, there is an exact sequence
\[
 0\longrightarrow\Br(k)\longrightarrow\bigoplus_v\Br(k_v)
 \longrightarrow\Q/\Z\longrightarrow0.
\]
The sum is over all primes of $k$, including the infinite primes. The first
map sends $[D]$ to $([D\otimes_k k_v])_v$, and the second sends
$(a_v)_v$ to $\sum_v\operatorname{inv}_v(a_v)$.
\end{theorem}

\begin{proof}[Source proof]

See CFT VIII, 4.2; Milne notes that this is no easier to prove than the main
theorem of abelian class field theory.

\end{proof}

\begin{theorem}[Brauer exponent]\label{thm:milne-brauer-exponent}\dcref{MI:ch2:2.8}
For a central division algebra $D$ over a number field $k$, the order of
$[D]$ in $\Br(k)$ equals $\sqrt{[D:k]}$ and the least common denominator of
its local invariants.
\end{theorem}

\begin{proof}[Source proof]

Since the order of an element of $\bigoplus_v\Q/\Z$ is the least common
denominator of its components, the second assertion follows directly from the
global Brauer exact sequence.  The first assertion follows from the
Grunwald--Wang theorem (CFT VIII, 2.4); Milne also refers to Reiner,
\emph{Maximal Orders}, 32.19.

\end{proof}

\begin{theorem}[Endomorphism algebra of a simple finite-field variety]\label{thm:milne-simple-finite-field}\dcref{MI:ch2:2.9}
Let $A/\F_q$ be simple, let $D=\End^0(A)$, and let $\pi$ be its Frobenius.
Then the centre of $D$ is $E=\Q[\pi]$. For a prime $v$ of $E$, put
$i_v=\operatorname{inv}_v(D)$. Then $|\pi|_v=q^{i_v}$, and, at a finite
place $v$ above $p$,
\[
\operatorname{inv}_v(D)=
\frac{\operatorname{ord}_v(\pi)}{\operatorname{ord}_v(q)}[E_v:\Q_p]
\quad\text{in }\Q/\Z;
\]
the invariant is $1/2$ at a real place and $0$ at a complex place or at a
finite place not dividing $p$.
Finally
$2\dim A=[D:E]^{1/2}[E:\Q]$.
\end{theorem}

\begin{proof}[Source proof]

This was proved by Tate (Inventiones Mathematicae, 1966), who did not publish
the proof of the assertion about the local invariants.  That part is given in
the account of Waterhouse and Milne, Proc. Symp. Pure Math., AMS XX.

\end{proof}

\begin{definition}[CM fields and CM types]\label{def:milne-cm-type}\dcref{MI:ch2:2}
A CM-field is a totally imaginary quadratic extension $E$ of a totally real
field $F$. Equivalently, it has an involution $\iota\ne 1$ whose action under
each embedding $E\hookrightarrow\C$ is complex conjugation. If $[E:\Q]=2g$,
a CM-type $\Phi$ is a choice of one embedding from each of the $g$ conjugate
pairs of embeddings $E\hookrightarrow\C$.
\end{definition}

\begin{definition}[Complex multiplication and good reduction]\label{def:milne-cm-good-reduction}\dcref{MI:ch2:2}
An abelian variety $A/\C$ of dimension $g$ has complex multiplication by a
CM-field $E$ if $E\hookrightarrow\End^0(A)$, $[E:\Q]=2g$, and a polarization
has Rosati involution restricting to the CM involution on $E$. Let $A$ be an
abelian variety over a nonarchimedean local field $k$, with valuation ring
$R$ and residue field $k_0$. Choose a projective embedding of $A$ and reduce
its defining equations modulo the maximal ideal of $R$, obtaining a variety
$A_0/k_0$. If the embedding can be chosen so that $A_0$ is nonsingular, then
$A_0$ is independent of all choices and is an abelian variety. In this case
$A$ is said to have good reduction and $A_0$ is its reduced variety. For an
abelian variety over a number field, good reduction at a finite prime means
good reduction after base change to the corresponding completion.
\end{definition}

\begin{proposition}[CM-type construction]\label{prop:milne-cm-type-construction}\dcref{MI:ch2:2}
Let $(E,\Phi)$ be a CM-type with $[E:\Q]=2g$. The quotient
\[
 A_\Phi=\C^g/\Phi(\mathcal O_E),\qquad
 \Phi(x)=(\phi(x))_{\phi\in\Phi},
\]
is a complex abelian variety of dimension $g$ with complex multiplication by
$E$.
\end{proposition}

\begin{proof}[Source proof]
The lattice $\Phi(\mathcal O_E)$ makes the quotient a complex torus. The
standard alternating form obtained from the trace pairing and the CM
involution is a Riemann form, so the torus is an abelian variety. Multiplication
by $x\in\mathcal O_E$ acts through the diagonal map $\Phi(x)$ and extends to
$E\subset\End^0(A_\Phi)$; the trace form gives the required Rosati involution.
\end{proof}

\begin{proposition}[CM Weil integer construction]\label{prop:milne-cm-weil-integer}\dcref{MI:ch2:2}
Let $(E,\Phi)$ be a CM-type, let $\mathfrak p$ be a prime of $E$ above the
rational prime $p$, and choose $h$ and $a\in\mathcal O_E$ with
$\mathfrak p^h=(a)$. For $n$ sufficiently divisible, set
\[
 \pi=\prod_{\phi\in\Phi}\phi(a^{2n}),
 \qquad q=p^{2nf(\mathfrak p/p)}.
\]
Then $\pi$ is a Weil $q$-integer, and its conjugacy class is independent of
the choice of the generator $a$.
\end{proposition}

\begin{proof}[Source proof]
Writing $\overline{\pi}$ for the product over the conjugate embeddings gives
\[
 \pi\overline{\pi}=\Norm_{E/\Q}(a^n)^2=q.
\]
The ideal norm calculation gives
$(\Norm_{E/\Q}a)=(p)^{f(\mathfrak p/p)}$, so the displayed integer is
$q=p^{2nf(\mathfrak p/p)}$. The same calculation for every conjugate of
$\pi$ proves the Weil absolute-value condition.

The unit theorem gives
$\operatorname{rank}(\mathcal O_E^\times)=g-1=\operatorname{rank}(\mathcal O_F^\times)$.
Choose $n$ so that $u^n\in\mathcal O_F^\times$ for every unit $u$ of $E$.
If $a$ is replaced by $au$, then, because the restrictions of the embeddings
in $\Phi$ are all the embeddings of $F$ into $\R$,
\[
 \prod_{\phi\in\Phi}\phi((au)^{2n})
 =\pi\,\Norm_{F/\Q}(u^n)^2=\pi.
\]
Thus the construction is independent of the choice of the generator.
\end{proof}

\begin{proposition}[CM reduction]\label{prop:milne-cm-reduction}\dcref{MI:ch2:2.10}
If $A/\C$ has complex multiplication by a CM-field $E$, then $A$ has a model over a number field and, after finite extension, good reduction at every prime.
\end{proposition}
\begin{proof}[Source proof omitted]
Milne writes ``Omitted'' and refers to the theory of complex multiplication.
\end{proof}
\begin{theorem}[Shimura--Taniyama]\label{thm:milne-shimura-taniyama}\dcref{MI:ch2:2.11}
Let $A/\F_q$ be obtained by reduction from an abelian variety of CM-type
$(E,\Phi)$. Then the Weil $q$-integer associated with $A$ is the one
constructed from a prime $\mathfrak p$ of $E$ and a generator of a power of
$\mathfrak p$, up to a root of unity.
\end{theorem}

\begin{proof}[Source proof]

This is the main theorem of Shimura and Taniyama,
\emph{Complex Multiplication of Abelian Varieties and its Applications to
Number Theory} (1961).  Milne refers to that monograph and, for a concise
modern exposition, to his note ``The fundamental theorem of complex
multiplication,'' arXiv:0705.3446.

\end{proof}
